\subsection{Milestone 3 -- Risultati dell'analisi controfattuale}

\subsubsection{Tabella dei dataset}

La Tabella~\ref{tab:whatif_datasets} riporta, per ciascun dataset, il numero di istanze,
i bug reali (ground truth da M1) e i bug predetti da BClassifierA.

\begin{table}[h]
\centering
\begin{tabular}{lrrr}
\toprule
\textbf{Dataset} & \textbf{Istanze} & \textbf{Actual Buggy} & \textbf{Predicted Buggy} \\
\midrule
A  (intero)              & 15.338 & 1.378 & 1.384 \\
B+ (NSmells $> 0$)       & 12.792 & 1.119 & 1.125 \\
B  (B+ con NSmells $= 0$) & 12.792 & 1.119 & 1.137 \\
C  (NSmells $= 0$)       &  2.546 &   259 &   259 \\
\bottomrule
\end{tabular}
\caption{Riepilogo dei dataset e predizioni di BClassifierA (RF + No FS + SMOTE addestrato su A).}
\label{tab:whatif_datasets}
\end{table}

I totali sono internamente consistenti: B+ + C = 12.792 + 2.546 = 15.338 = A, e B ha lo
stesso numero di istanze di B+ per costruzione.

\subsubsection{Metriche di prevenzione}

\begin{table}[h]
\centering
\begin{tabular}{lr}
\toprule
\textbf{Metrica} & \textbf{Valore} \\
\midrule
Prevented\_Total           &  63    \\
Prevented\_Proportion      &  4,57\% \\
Prevented\_Of\_Predictable &  5,60\% \\
\bottomrule
\end{tabular}
\caption{Metriche di prevenzione: impatto dell'azzeramento di NSmells sulle predizioni.}
\label{tab:whatif_prevention}
\end{table}

Su 1.125 classi predette buggy in B+, solo \textbf{63} (5,60\%) smetterebbero di essere
predette buggy se i code smell fossero zero. In proporzione sul totale dei bug reali in A,
la prevenzione è del 4,57\%.

\subsubsection{Un risultato contro-intuitivo: B predice più bug di B+}

Un dato inatteso emerge dal confronto diretto: B ha \textbf{1.137} predizioni buggy contro
le \textbf{1.125} di B+ ($+12$ netto). L'azzeramento degli smell non riduce le predizioni
in modo uniforme: 63 istanze passano da BUGGY a CLEAN (prevenute), ma \textbf{75 istanze}
percorrono la direzione opposta, da CLEAN a BUGGY. L'effetto netto dell'eliminazione degli
smell è un lieve \textit{aumento} del numero di bug predetti, a conferma del ruolo marginale
di NSmells nel modello.

\subsubsection{Confronto con il paper di riferimento}

Il paper presentato a lezione (\textit{``What If I Had No Smells?''}, Falessi) riporta una
riduzione del \textbf{42\%} dei file difettosi predetti passando da B+ a B, contro il
nostro 5,60\%. La differenza è sostanziale e richiede una spiegazione.

La metodologia è la stessa: A, B+, B, C costruiti con lo stesso criterio, modello addestrato
su A e applicato ai quattro dataset. Le differenze risiedono nel contesto sperimentale:

\begin{itemize}

  \item \textbf{Strumento di rilevazione smell e composizione dei ruleset.}
  Il paper usa \textbf{SonarQube} con 1.297 tipi di smell, che include veri smell
  architetturali (God Class, Long Method, Feature Envy, ecc.). Il presente studio usa
  \textbf{PMD} con 4 ruleset: \texttt{design}, \texttt{errorprone}, \texttt{bestpractices}
  e \texttt{codestyle}. I primi tre catturano pattern di qualità strutturale (God Class,
  Cyclomatic Complexity, pattern error-prone); il quarto --- \texttt{codestyle} ---
  rileva violazioni di stile puro: naming convention, stile delle parentesi graffe,
  import superflui. Le violazioni di codestyle crescono \textit{linearmente} con la
  dimensione del file: un file da 500 righe ha più opportunità di violare una naming
  convention di uno da 50 righe. Questo introduce una correlazione artificiale tra
  NSmells e LOC, che a sua volta è correlata con le metriche di processo cumulative
  (LOC\_touched, NR, Churn). In altri termini, \texttt{codestyle} riduce l'indipendenza
  di NSmells dalle metriche di processo, penalizzando la sua utilità come feature
  indipendente nel modello. Al contrario, le violazioni di \texttt{design} (es.\ God
  Class) non crescono linearmente con la dimensione e conservano un segnale più autonomo.

  \item \textbf{Contributo marginale di NSmells nel modello multivariato.}
  Nella slide di correlazione del paper, NSmells presenta correlazione quasi nulla con
  LOC\_touched ($-0{,}01$), NR ($0{,}02$) e NFix ($0{,}02$): nel loro dataset i code smell
  sono \textit{indipendenti} dalle metriche di processo, e il loro contributo al modello
  è quindi aggiuntivo rispetto a quello delle altre feature.

  È importante distinguere due concetti diversi. La \textit{correlazione individuale} di
  NSmells con isBuggy (misurata da InfoGain o Spearman) è positiva e non nulla anche nel
  presente studio: per questo la feature selection con soglia $= 0$ non elimina NSmells,
  e il classificatore migliore risulta comunque quello senza FS (FS è un no-op quando
  tutte le feature hanno InfoGain $> 0$). Il \textit{contributo marginale} di NSmells
  nel modello --- cioè quanto cambia la predizione se si varia NSmells \textit{tenendo
  fisse} le altre 18 feature --- è invece quasi zero: RandomForest ha già spiegato la
  quasi totalità della varianza di isBuggy tramite le metriche di processo; NSmells non
  aggiunge informazione indipendente. È questo contributo marginale, non la correlazione
  individuale, che il what-if misura e che risulta marginale.

  \item \textbf{Contesto industriale vs.\ open source.}
  Il paper usa dati di un'azienda industriale (Keymind): 135 file totali, processo di
  sviluppo più controllato. Il presente studio usa Apache OpenJPA: 15.338 snapshot, sviluppo
  open source distribuito ad alto turnover. In contesti industriali disciplinati le metriche
  di churn sono più basse e gli smell diventano relativamente più informativi; in grandi
  progetti open source il churn domina la predizione (NFix ha correlazione $0{,}51^{*}$ con
  la bugginess già nel paper del prof).

  \item \textbf{Sfasamento temporale della metrica.}
  Il paper misura NSmells \textit{just before the first revision} della finestra di release,
  cioè allo stato iniziale preciso della release. Il presente studio usa
  \texttt{Previous\_Release\_Code\_Smells}, ossia i smell misurati sulla release
  \textit{precedente}: uno sfasamento che può diluire il segnale se il codice è cambiato
  significativamente tra una release e l'altra.

\end{itemize}

\subsubsection{Possibili miglioramenti della metrica NSmells}

L'analisi evidenzia due interventi che potrebbero avvicinare i risultati a quelli del paper:

\begin{itemize}
  \item \textbf{Rimozione di \texttt{codestyle.xml} e aggiunta di ruleset più smell-oriented.}
  Eliminare il ruleset \texttt{codestyle} ridurrebbe il rumore correlato alla dimensione;
  aggiungere \texttt{performance.xml} e \texttt{security.xml} introdurrebbe violazioni di
  qualità più indipendenti dalle metriche di processo (pattern di performance e sicurezza non
  crescono linearmente con il numero di righe). Il risultato sarebbe una metrica NSmells più
  focalizzata su smell strutturali e meno confusa con lo stile.

  \item \textbf{Sostituzione di PMD con SonarQube.}
  SonarQube rileva 1.297 tipi di smell, inclusi quelli architetturali (God Class, Feature
  Envy, Divergent Change, ecc.) che nella letteratura hanno la correlazione più alta con i
  difetti. Questa sostituzione richiederebbe un refactoring significativo della pipeline
  (server SonarQube, analisi per release, parsing API) ma sarebbe l'unico intervento capace
  di colmare strutturalmente il divario con il 42\% riportato nel paper.
\end{itemize}

Entrambi gli interventi sono stati valutati ma non implementati nel presente studio,
che mantiene la scelta originale di PMD per coerenza con la pipeline di Milestone 1.

\paragraph{Risposta alle domande dell'assignment.}
\textit{``BClassifier è accurato?''} Sì: Precision $= 0{,}63$, Recall $= 0{,}66$,
AUC $= 0{,}94$, Kappa $= 0{,}61$ da M2.
\textit{``Quante classi buggy avrebbero potuto essere prevenute con zero smell?''}
63 in totale (4,57\% dei bug reali), pari al 5,60\% dei bug predetti in B+.
L'impatto è limitato: nel contesto di Apache OpenJPA con metriche PMD, il principale
driver di bugginess sono le metriche di processo cumulative, non la qualità strutturale
del codice misurata tramite PMD.

\subsubsection{Analisi di ridondanza (M3.1)}

Per spiegare empiricamente il basso impatto di NSmells, sono state condotte due analisi
supplementari sul dataset A.

\paragraph{Correlazione di Spearman tra NSmells e le altre feature.}

La Tabella~\ref{tab:whatif_correlation} riporta la correlazione di Spearman $\rho$ tra
\texttt{Previous\_Release\_Code\_Smells} e ciascuna delle altre 18 feature, ordinate per
valore assoluto decrescente.

\begin{table}[h]
\centering
\begin{tabular}{lr}
\toprule
\textbf{Feature} & $\boldsymbol{\rho}$ \\
\midrule
LOC               & $0{,}7821$ \\
Max\_Churn        & $0{,}7123$ \\
LOC\_Touched      & $0{,}6909$ \\
Churn             & $0{,}6909$ \\
Weighted\_Age     & $0{,}6900$ \\
Max\_LOC\_Added   & $0{,}6632$ \\
LOC\_Added        & $0{,}6618$ \\
Avg\_Churn        & $0{,}5621$ \\
Nauth             & $0{,}5312$ \\
NR                & $0{,}5276$ \\
Avg\_LOC\_Added   & $0{,}4700$ \\
Age               & $0{,}2894$ \\
Change\_Set\_Size & $0{,}2518$ \\
Avg\_Change\_Set  & $-0{,}2416$ \\
Nfix              & $0{,}1666$ \\
Max\_Change\_Set  & $0{,}1606$ \\
SEXP              & $0{,}1302$ \\
EXP               & $0{,}0381$ \\
\bottomrule
\end{tabular}
\caption{Correlazione di Spearman tra \texttt{Previous\_Release\_Code\_Smells} e le altre
         18 feature del dataset A (Apache OpenJPA, 15.338 istanze).}
\label{tab:whatif_correlation}
\end{table}

Sei feature su 18 presentano $\rho > 0{,}60$ con NSmells: LOC ($0{,}78$), Max\_Churn
($0{,}71$), LOC\_Touched ($0{,}69$), Churn ($0{,}69$), Weighted\_Age ($0{,}69$),
Max\_LOC\_Added ($0{,}66$). Il risultato è coerente con il meccanismo descritto nella
sezione precedente: le violazioni del ruleset \texttt{codestyle} di PMD crescono
proporzionalmente alla dimensione del file, introducendo una correlazione artificiale tra
NSmells e tutte le metriche che catturano dimensione e storia delle modifiche. Ciò riduce
l'indipendenza informativa di NSmells rispetto alle metriche di processo, spiegando
perché RandomForest non la utilizza attivamente.

Le uniche feature con correlazione bassa rispetto a NSmells sono EXP ($0{,}04$) e SEXP
($0{,}13$), che misurano l'esperienza dell'autore e sono intrinsecamente indipendenti dalla
dimensione del file. Tuttavia, anche queste feature hanno correlazione individuale con
isBuggy bassa, per cui non compensano la debolezza di NSmells nel modello.

\paragraph{Studio di ablazione: RF+SMOTE con e senza NSmells.}

La Tabella~\ref{tab:whatif_ablation} confronta le prestazioni di RF + No FS + SMOTE in
10$\times$10-fold CV sul dataset A completo (19 feature) e sul dataset A privo di NSmells
(18 feature).

\begin{table}[h]
\centering
\begin{tabular}{lrrrrrr}
\toprule
\textbf{Configurazione} & \textbf{Feature} & \textbf{P} & \textbf{R} & \textbf{AUC} & \textbf{Kappa} & \textbf{NPofB20} \\
\midrule
RF + SMOTE (tutte le feature) & 19 & $0{,}625$ & $0{,}664$ & $0{,}942$ & $0{,}608$ & $0{,}880$ \\
RF + SMOTE (senza NSmells)    & 18 & $0{,}629$ & $0{,}664$ & $0{,}941$ & $0{,}610$ & $0{,}880$ \\
\midrule
$\Delta$ (senza $-$ con)      & $-1$ & $+0{,}004$ & $\approx 0$ & $\approx 0$ & $+0{,}002$ & $\approx 0$ \\
\bottomrule
\end{tabular}
\caption{Studio di ablazione: prestazioni di RF + No FS + SMOTE in 10$\times$10-fold CV
         con e senza la feature \texttt{Previous\_Release\_Code\_Smells}.}
\label{tab:whatif_ablation}
\end{table}

Il $\Delta$ è essenzialmente zero su tutte le metriche. In modo controintuitivo, rimuovere
NSmells produce un lieve \textit{miglioramento} di Precision ($+0{,}004$) e Kappa
($+0{,}002$): quando una feature è ridondante e collineare con altre, il modello la può
usare in modo rumoroso, suddividendo lo spazio delle decisioni in modo leggermente meno
pulito. Rimuoverla riduce questo rumore marginale. L'effetto è minimo ma in direzione
opposta a quella attesa se NSmells avesse valore indipendente.

La configurazione a 19 feature produce Kappa $= 0{,}608$ e AUC $= 0{,}942$, sostanzialmente
identici ai valori di M2 (Kappa $= 0{,}61$, AUC $= 0{,}94$): la CV è riproducibile e le due
esecuzioni sono coerenti.

\paragraph{Sintesi delle evidenze M3.1.}
Le due analisi si rinforzano a vicenda: NSmells correla fortemente con le metriche di
processo ($\rho$ fino a $0{,}78$ con LOC) \textit{e} rimuoverla non degrada il modello
($\Delta \approx 0$). L'insieme di queste evidenze dimostra empiricamente che il basso
impatto del what-if di M3 (5,60\%) non dipende da un modello poco accurato, ma dalla
ridondanza informativa di NSmells rispetto alle metriche di processo già presenti nel
dataset.
